\section{Correct signed warm-start SOR analysis}
\label{sec:sor-correction}

A central correction from the discussion is the distinction between two
parameter ranges:
\begin{equation}
    \boxed{
    0<\omega<2
    \text{ is the classical global SOR convergence range,}}
    \label{eq:global-sor-range}
\end{equation}
whereas a graph-uniform one-coordinate weighted-\(\ell_1\) contraction holds
only on a smaller interval.  Conflating these statements leads to an invalid
local work proof for the optimal over-relaxed parameter.

\begin{proposition}[Correct signed weighted-\(\ell_1\) bound]
\label{prop:signed-sor-l1}
Let \(q=D^{1/2}g(x)\), \(\chi=(1-\alpha)/(1+\alpha)\), and apply the relaxed
coordinate update \cref{eq:general-coordinate-update}.  Then
\begin{equation}
    q^+
    =q-\omega q_u e_u
      +\chi\omega q_uAD^{-1}e_u.
    \label{eq:signed-sor-scaled-update}
\end{equation}
For every signed vector \(q\),
\begin{equation}
    \norm{q^+}_1
    \leq
    \norm{q}_1-\delta_\alpha(\omega)\abs{q_u},
    \label{eq:signed-sor-bound}
\end{equation}
where
\begin{equation}
    \delta_\alpha(\omega)
    :=1-\abs{1-\omega}-\chi\omega
    =
    \begin{cases}
      \dfrac{2\alpha\omega}{1+\alpha},
      &0<\omega\leq1,\\[2mm]
      \dfrac{2(1+\alpha-\omega)}{1+\alpha},
      &1\leq\omega<2.
    \end{cases}
    \label{eq:delta-alpha}
\end{equation}
Hence strict uniform weighted-\(\ell_1\) decrease is guaranteed precisely for
\begin{equation}
    0<\omega<1+\alpha.
    \label{eq:l1-safe-range}
\end{equation}
\end{proposition}

\begin{proof}
Only coordinate \(u\) and its neighbors change.  Since the graph is simple,
\(q_u^+=(1-\omega)q_u\), while
\begin{equation}
    q_v^+
    =q_v+\chi\omega\frac{q_u}{d_u}
    \qquad(v\sim u).
    \label{eq:sor-neighbor-update}
\end{equation}
Using \(\norm{AD^{-1}e_u}_1=1\),
\begin{align*}
  \norm{q^+}_1
  &\leq
  \norm{q}_1-\abs{q_u}
  +\abs{1-\omega}\abs{q_u}
  +\chi\omega\abs{q_u}\\
  &=\norm{q}_1-
    \bigl(1-\abs{1-\omega}-\chi\omega\bigr)\abs{q_u}.
\end{align*}
Expanding the absolute value on the intervals \((0,1]\) and \([1,2)\)
gives \cref{eq:delta-alpha}.  The coefficient is positive exactly when
\(\omega<1+\alpha\).
\end{proof}

\begin{remark}[Why the earlier derivation failed]
The invalid step was to replace
\(
\abs{(1-\omega)q_u}
\)
by
\(
(1-\omega)\abs{q_u}
\)
for \(\omega>1\).  The correct value is
\((\omega-1)\abs{q_u}\).  The formula valid on \(0<\omega\leq1\) therefore
cannot be extended unchanged to optimal over-relaxation.
\end{remark}

\begin{proposition}[Sharpness of the safe range]
\label{prop:sor-l1-sharp}
For every \(\omega>1+\alpha\), there exists a signed warm-start residual for
which one SOR coordinate update increases \(\norm{q}_1\).
\end{proposition}

\begin{proof}
Take \(q\) supported only on the processed vertex \(u\).  Then all triangle
inequalities in \cref{prop:signed-sor-l1} are equalities, and
\begin{equation}
    \frac{\norm{q^+}_1}{\norm{q}_1}
    =\abs{1-\omega}+\chi\omega.
    \label{eq:sor-supported-ratio}
\end{equation}
For \(\omega>1\), this equals
\(
2\omega/(1+\alpha)-1
\), which exceeds one exactly when \(\omega>1+\alpha\).
\end{proof}

The optimal global parameter from \cref{eq:optimal-sor} lies in the uniform
\(\ell_1\)-safe range only when
\begin{equation}
    \omega_\star\leq1+\alpha
    \quad\Longleftrightarrow\quad
    \alpha\geq(\sqrt2-1)^2=3-2\sqrt2\approx0.1716.
    \label{eq:omega-star-safe-condition}
\end{equation}
For the small \(\alpha\) regime where acceleration is most valuable,
\(\omega_\star\) is therefore outside the graph-uniform signed
weighted-\(\ell_1\) contraction interval.  Its speed comes from multi-step
spectral effects and cancellation, not a larger guaranteed one-coordinate
mass decrease.

\subsection{An exact local safeguard}

The failure of a universal \(\ell_1\) theorem does not prevent using
\(\omega_\star\) opportunistically.  Because only \(u\) and \(N(u)\) change,
the exact local mass decrease of a candidate parameter \(\omega\) is
\begin{align}
    \Delta_u(\omega)
    &:={}
    \abs{q_u}+\sum_{v\sim u}\abs{q_v}
    -\abs{(1-\omega)q_u}
    -\sum_{v\sim u}
       \abs{q_v+\chi\omega q_u/d_u}.
    \label{eq:exact-local-mass-decrease}
\end{align}
Computing \(\Delta_u(\omega)\) costs \(\mathcal O(d_u)\), the same order as
executing the push.

\begin{definition}[Safeguarded optimal LocSOR]
\label{def:safeguarded-sor}
Fix \(\sigma\in(0,1]\) and let
\begin{equation}
    \delta_{\rm GS}:=\frac{2\alpha}{1+\alpha}.
    \label{eq:gs-delta}
\end{equation}
At an active coordinate \(u\), accept the \(\omega_\star\) candidate only if
\begin{equation}
    \Delta_u(\omega_\star)
    \geq\sigma\delta_{\rm GS}\abs{q_u}.
    \label{eq:safeguard-test}
\end{equation}
Otherwise execute the \(\omega=1\) update.
\end{definition}

\begin{theorem}[Safeguard guarantee]
\label{thm:safeguarded-sor}
Every committed update of safeguarded optimal LocSOR satisfies
\begin{equation}
    M(x^+)\leq M(x)
      -\sigma\frac{2\alpha}{1+\alpha}\abs{q_u}.
    \label{eq:safeguard-decrease}
\end{equation}
Consequently, from any handoff point \(x_0\),
\begin{equation}
    \Work_{\rm tail}^{\rm safe}
    \leq
    \frac{1+\alpha}{2\sigma\alpha^2\eps}M(x_0).
    \label{eq:safeguard-tail-work}
\end{equation}
Under the mass handoff
\(M(x_0)\leq\theta\alpha^{3/2}\), this becomes
\begin{equation}
    \Work_{\rm tail}^{\rm safe}
    \leq
    \frac{\theta(1+\alpha)}{2\sigma\sqrt\alpha\eps}.
    \label{eq:safeguard-accelerated}
\end{equation}
\end{theorem}

\begin{proof}
If the \(\omega_\star\) candidate is accepted, the claim is exactly the test
\cref{eq:safeguard-test}.  If it is rejected, the fallback \(\omega=1\)
satisfies the stronger decrease with \(\sigma=1\) by
\cref{lem:gs-mass-decrease}.  An active coordinate obeys
\(\abs{q_u}\geq\alpha\eps d_u\).  Sum the resulting decrease over all
committed updates.
\end{proof}

\begin{remark}[Recommended theoretical hierarchy]
The cleanest correctness and work theorem uses \(\omega=1\).  The safeguarded
method preserves the same asymptotic guarantee while testing optimal
relaxation at no additional asymptotic edge cost.  Plain \(\omega_\star\)
remains a valuable empirical baseline and is globally convergent, but the
weighted-\(\ell_1\) proof must not be assigned to it.
\end{remark}
